% 56e Olympiade Internationale de Mathématiques — Chiang Mai, Thaïlande.
% Énoncés officiels en français (imo-official.org). Chaque problème vaut 7 points.

\begin{center}
  \emph{Premier jour --- vendredi 10 juillet 2015 --- Durée : 4 heures 30}
\end{center}

\subsection*{Problème 1}
On dit qu'un ensemble fini $\mathcal{S}$ de points du plan est \emph{équilibré} si, pour tous points $A$ et $B$
de $\mathcal{S}$ distincts, il existe un point $C$ de $\mathcal{S}$ tel que $AC = BC$. On dit que $\mathcal{S}$
est \emph{excentrique} si, pour tous points $A$, $B$ et $C$ de $\mathcal{S}$ distincts, il n'existe pas de
point $P$ de $\mathcal{S}$ tel que $PA = PB = PC$.
\begin{enumerate}
  \item[(a)] Prouver que pour tout entier $n \geqslant 3$, il existe un ensemble équilibré contenant exactement
  $n$ points.
  \item[(b)] Déterminer tous les entiers $n \geqslant 3$ pour lesquels il existe un ensemble équilibré et
  excentrique contenant exactement $n$ points.
\end{enumerate}

\subsection*{Problème 2}
Déterminer tous les triplets $(a, b, c)$ d'entiers strictement positifs pour lesquels chacun des nombres
\[
ab - c, \quad bc - a, \quad ca - b
\]
est une puissance de $2$.

(\emph{Une puissance de $2$ est un entier de la forme $2^n$, où $n$ est un entier positif ou nul.})

\subsection*{Problème 3}
Soit $ABC$ un triangle dont tous les angles sont aigus, avec $AB > AC$. Soit $\Gamma$ son cercle circonscrit,
$H$ son orthocentre et $F$ le pied de sa hauteur issue de $A$. On désigne par $M$ le milieu du segment $[BC]$.
Soit $Q$ le point de $\Gamma$ tel que $\widehat{HQA} = 90^\circ$ et soit $K$ le point de $\Gamma$ tel que
$\widehat{HKQ} = 90^\circ$. On suppose que les points $A$, $B$, $C$, $K$ et $Q$ sont tous distincts et dans cet
ordre sur $\Gamma$.

Prouver que le cercle circonscrit au triangle $KQH$ est tangent au cercle circonscrit au triangle $FKM$.

\bigskip
\begin{center}
  \emph{Second jour --- samedi 11 juillet 2015 --- Durée : 4 heures 30}
\end{center}

\subsection*{Problème 4}
Soit $ABC$ un triangle de cercle circonscrit $\Omega$, et soit $O$ le centre de $\Omega$. Un cercle $\Gamma$ de
centre $A$ rencontre le segment $[BC]$ aux points $D$ et $E$, de sorte que $B$, $D$, $E$ et $C$ sont distincts
et dans cet ordre sur la droite $(BC)$. On note $F$ et $G$ les points d'intersection de $\Gamma$ et $\Omega$,
de sorte que $A$, $F$, $B$, $C$ et $G$ sont dans cet ordre sur $\Omega$. Soit $K$ le second point
d'intersection du cercle circonscrit au triangle $BDF$ avec le segment $[AB]$. Soit $L$ le second point
d'intersection du cercle circonscrit au triangle $CGE$ avec le segment $[CA]$.

On suppose que les droites $(FK)$ et $(GL)$ ne sont pas confondues et qu'elles se rencontrent au point $X$.
Prouver que $X$ appartient à la droite $(AO)$.

\subsection*{Problème 5}
Soit $\mathbb{R}$ l'ensemble des nombres réels. Déterminer toutes les fonctions
$f : \mathbb{R} \to \mathbb{R}$ qui vérifient l'équation
\[
f\bigl(x + f(x + y)\bigr) + f(xy) \ = \ x + f(x + y) + yf(x)
\]
pour tous réels $x$ et $y$.

\subsection*{Problème 6}
La suite $a_1, a_2, \ldots$ d'entiers vérifie les conditions :
\begin{enumerate}
  \item[(i)] $1 \leqslant a_j \leqslant 2015$ pour tout $j \geqslant 1$,
  \item[(ii)] $k + a_k \neq \ell + a_\ell$ pour tous $1 \leqslant k < \ell$.
\end{enumerate}
Prouver qu'il existe deux entiers strictement positifs $b$ et $N$ pour lesquels
\[
\left|\sum_{j=m+1}^{n} (a_j - b)\right| \leqslant 1007^2
\]
pour tous les entiers $m$ et $n$ tels que $n > m \geqslant N$.
